\documentclass[12pt,a4paper]{article}
\usepackage[scheme=chinese]{ctex}
\usepackage{geometry}
\usepackage{graphicx}
\usepackage{booktabs}
\usepackage{amsmath,amssymb}
\usepackage{listings}
\usepackage{xcolor}
\usepackage{float}
\usepackage{longtable}
\usepackage{array}
\usepackage{multirow}
\usepackage{hyperref}
\usepackage{caption}
\usepackage{subcaption}
\usepackage{enumitem}
\usepackage{fancyhdr}
\usepackage{titlesec}
\usepackage{setspace}

\geometry{left=2.5cm,right=2.5cm,top=2.5cm,bottom=2.5cm}

\setlength{\parindent}{2em}
\linespread{1.25}

\titleformat{\section}{\centering\heiti\fontsize{16pt}{\baselineskip}\selectfont}{第\chinese{section}章}{1em}{}
\titleformat{\subsection}{\heiti\fontsize{14pt}{\baselineskip}\selectfont}{\thesubsection}{1em}{}
\titleformat{\subsubsection}{\heiti\fontsize{13pt}{\baselineskip}\selectfont}{\thesubsubsection}{1em}{}

\titlespacing{\section}{0pt}{24pt}{18pt}
\titlespacing{\subsection}{0pt}{24pt}{6pt}
\titlespacing{\subsubsection}{0pt}{12pt}{6pt}

\lstset{
    language=Python,
    basicstyle=\ttfamily\fontsize{9pt}{11pt}\selectfont,
    keywordstyle=\color{blue}\bfseries,
    commentstyle=\color{green!60!black},
    stringstyle=\color{red!70!black},
    numbers=left,
    numberstyle=\tiny\color{gray},
    numbersep=5pt,
    frame=single,
    framesep=3pt,
    breaklines=true,
    breakatwhitespace=true,
    showstringspaces=false,
    tabsize=4,
    captionpos=b,
    backgroundcolor=\color{gray!5},
    xleftmargin=1em,
    xrightmargin=1em
}

\begin{document}

\begin{center}
{\heiti\fontsize{18pt}{22pt}\selectfont\bfseries 强化学习第三次作业}
\end{center}

\section*{摘要}
\addcontentsline{toc}{section}{摘要}

{\songti\fontsize{12pt}{15pt}\selectfont 本实验针对CartPole环境，实现了策略迭代算法和值迭代算法，用于求解最优控制策略。由于CartPole是连续状态空间环境，本实验首先对状态空间进行离散化处理，将其转化为有限状态的马尔可夫决策过程（MDP）。然后实现了两种策略评估方法：解析法和数值法。在此基础上，分别实现了策略迭代和值迭代两种动态规划算法，并对两种算法的收敛性、计算效率进行了详细比较。实验结果表明，两种算法均能收敛到有效的控制策略，策略迭代平均奖励达到165.44，值迭代达到179.71，相比随机策略基线（23.89）分别提升592.5\%和652.2\%。本实验加深了对强化学习中动态规划方法的理解，为后续学习更复杂的强化学习算法奠定了基础。}

\vspace{0.5em}
{\heiti\fontsize{12pt}{15pt}\selectfont 关键词：}{\songti\fontsize{12pt}{15pt}\selectfont CartPole；状态离散化；策略迭代；值迭代；动态规划}

\section*{Abstract}
\addcontentsline{toc}{section}{Abstract}

{\fontsize{12pt}{15pt}\selectfont This experiment implements policy iteration and value iteration algorithms for the CartPole environment to find optimal control policies. Since CartPole is a continuous state space environment, this experiment first discretizes the state space to transform it into a finite-state Markov Decision Process (MDP). Two policy evaluation methods are then implemented: analytical method (solving linear equations) and numerical method (iterative update). Based on these, both policy iteration and value iteration algorithms are implemented, and their convergence and computational efficiency are compared in detail. Experimental results show that both algorithms converge to effective control policies, with policy iteration achieving an average reward of 165.44 and value iteration achieving 179.71, representing improvements of 592.5\% and 652.2\% respectively compared to the random policy baseline (23.89). This experiment deepens the understanding of dynamic programming methods in reinforcement learning and lays the foundation for learning more complex reinforcement learning algorithms.}

\vspace{0.5em}
{\heiti\fontsize{12pt}{15pt}\selectfont Key Words:} {\fontsize{12pt}{15pt}\selectfont CartPole; State Discretization; Policy Iteration; Value Iteration; Dynamic Programming}

\newpage

\section{引言}

\subsection{研究背景}

强化学习是机器学习的一个重要分支，其核心思想是智能体通过与环境交互，学习最优行为策略以最大化累积奖励。CartPole是强化学习领域的经典控制问题，智能体需要通过左右移动小车来保持杆子平衡。

马尔可夫决策过程（MDP）是强化学习的数学基础，为描述序贯决策问题提供了理论框架。一个MDP可以用五元组$(S, A, P, R, \gamma)$表示，其中$S$为状态空间，$A$为动作空间，$P$为状态转移概率，$R$为奖励函数，$\gamma$为折扣因子。

动态规划是求解MDP最优策略的经典方法，主要包括策略迭代和值迭代两种算法。本实验旨在通过实现这两种算法，深入理解其原理和特点。

\subsection{实验目的}

\begin{enumerate}[itemsep=0pt]
    \item 掌握MDP的基本概念和建模方法
    \item 理解并实现策略评估的解析法和数值法
    \item 实现策略迭代和值迭代算法
    \item 比较两种算法的收敛性和计算效率
    \item 将理论算法应用于实际控制问题
\end{enumerate}

\section{实验环境与MDP建模}

\subsection{CartPole环境介绍}

CartPole是一个经典的控制问题，环境描述如下：

\textbf{状态空间}：
\begin{itemize}[itemsep=0pt]
    \item 小车位置（cart\_position）：$[-2.4, 2.4]$
    \item 小车速度（cart\_velocity）：$(-\infty, +\infty)$
    \item 杆子角度（pole\_angle）：$[-0.2095, 0.2095]$ rad（约$\pm 12°$）
    \item 杆子角速度（pole\_angular\_velocity）：$(-\infty, +\infty)$
\end{itemize}

\textbf{动作空间}：
\begin{itemize}[itemsep=0pt]
    \item 动作0：向左施力
    \item 动作1：向右施力
\end{itemize}

\textbf{奖励函数}：每维持一步获得奖励+1，杆子倒下或小车出界则终止。

\subsection{状态离散化}

由于CartPole是连续状态空间，需要将其离散化为有限状态空间才能应用动态规划。本实验采用均匀分箱（Uniform Binning）方法进行离散化。

\begin{table}[H]
\centering
\caption{状态离散化参数设置}
\label{tab:discretization}
\begin{tabular}{cccc}
\toprule
\textbf{状态分量} & \textbf{原始范围} & \textbf{离散化范围} & \textbf{分箱数} \\
\midrule
小车位置 & $[-4.8, 4.8]$ & $[-2.4, 2.4]$ & 8 \\
小车速度 & $(-\infty, +\infty)$ & $[-3.0, 3.0]$ & 8 \\
杆子角度 & $[-0.418, 0.418]$ & $[-0.2, 0.2]$ & 8 \\
杆子角速度 & $(-\infty, +\infty)$ & $[-3.0, 3.0]$ & 8 \\
\bottomrule
\end{tabular}
\end{table}

离散化后，总状态数为$8^4 = 4096$个。离散化函数实现如下：

\begin{lstlisting}[caption={状态离散化函数}]
def discretize_state(self, state):
    cart_pos, cart_vel, pole_angle, pole_ang_vel = state
    pos_idx     = int(np.clip(np.digitize(cart_pos, 
                     self.position_bins), 0, self.n_bins - 1))
    vel_idx     = int(np.clip(np.digitize(cart_vel, 
                     self.velocity_bins), 0, self.n_bins - 1))
    angle_idx   = int(np.clip(np.digitize(pole_angle, 
                     self.angle_bins), 0, self.n_bins - 1))
    ang_vel_idx = int(np.clip(np.digitize(pole_ang_vel, 
                     self.angular_velocity_bins), 0, self.n_bins - 1))
    return (pos_idx * self.n_bins ** 3 +
            vel_idx * self.n_bins ** 2 +
            angle_idx * self.n_bins +
            ang_vel_idx)
\end{lstlisting}

\subsection{转移矩阵构建}

本实验采用从真实环境采样的方式构建转移矩阵。通过在CartPole环境中执行随机动作，记录状态转移信息，统计得到转移概率。

\begin{lstlisting}[caption={从真实环境采样构建转移矩阵}]
def _build_transition_matrix_from_env(self, n_samples):
    env = gym.make('CartPole-v1')
    trans_data = [[{} for _ in range(self.n_actions)] 
                  for _ in range(self.n_states)]
    
    while total_samples < n_samples:
        state, _ = env.reset()
        while not (terminated or truncated):
            s = self.discretize_state(state)
            action = env.action_space.sample()
            next_state, reward, terminated, truncated, _ = env.step(action)
            next_s = self.discretize_state(next_state)
            # 记录转移数据
            entry = trans_data[s][action]
            if next_s not in entry:
                entry[next_s] = {'count': 0, 'reward_sum': 0.0, 'done': False}
            entry[next_s]['count'] += 1
            entry[next_s]['reward_sum'] += reward
            if terminated:
                entry[next_s]['done'] = True
\end{lstlisting}

\subsection{MDP模型总结}

最终构建的MDP模型参数如表\ref{tab:mdp_params}所示。

\begin{table}[H]
\centering
\caption{MDP模型参数}
\label{tab:mdp_params}
\begin{tabular}{cc}
\toprule
\textbf{参数} & \textbf{值} \\
\midrule
状态数 $|S|$ & 4096 \\
动作数 $|A|$ & 2 \\
折扣因子 $\gamma$ & 0.99 \\
采样次数 & 50000 \\
\bottomrule
\end{tabular}
\end{table}

\section{算法原理与实现}

\subsection{策略评估}

策略评估是计算给定策略$\pi$的值函数$V^\pi$的过程。值函数定义为：
\begin{equation}
V^\pi(s) = \mathbb{E}_\pi\left[\sum_{t=0}^{\infty} \gamma^t r_t \mid s_0 = s\right]
\end{equation}

\subsubsection{解析法}

解析法通过求解线性方程组得到精确的值函数。根据贝尔曼方程：
\begin{equation}
V^\pi = R^\pi + \gamma P^\pi V^\pi
\end{equation}

整理得：
\begin{equation}
(I - \gamma P^\pi)V^\pi = R^\pi
\end{equation}

其中$P^\pi$是策略$\pi$下的状态转移矩阵，$R^\pi$是策略$\pi$下的期望奖励向量。

\begin{lstlisting}[caption={解析法策略评估}]
def evaluate_analytical(self, policy):
    n = self.mdp.n_states
    gamma = self.mdp.gamma
    
    # 计算策略下的期望奖励 R_pi
    R_pi = np.einsum('sa,sa->s', policy, self.mdp.R_mat)
    
    # 计算策略下的状态转移矩阵 P_pi
    P_pi = np.zeros((n, n))
    for a in range(na):
        for t in range(max_t):
            prob = self.mdp.T_prob[:, a, t]
            next_s = self.mdp.T_next[:, a, t]
            not_done = 1.0 - self.mdp.T_done[:, a, t]
            weight = policy[:, a] * prob * not_done
            np.add.at(P_pi, (idx, next_s), weight)
    
    # 求解线性方程组
    A = np.eye(n) - gamma * P_pi
    V = np.linalg.solve(A, R_pi)
    return V
\end{lstlisting}

\subsubsection{数值法}

数值法通过迭代更新逼近值函数：
\begin{equation}
V_{k+1}(s) = \sum_a \pi(a|s) \sum_{s'} P(s'|s,a)[r + \gamma V_k(s') \cdot (1-\text{done})]
\end{equation}

\begin{lstlisting}[caption={数值法策略评估}]
def evaluate_iterative(self, policy, tol=1e-6, max_iter=5000):
    V = np.zeros(self.mdp.n_states)
    for iteration in range(max_iter):
        Q = self.mdp.compute_q_values(V)
        V_new = np.einsum('sa,sa->s', policy, Q)
        delta = np.max(np.abs(V_new - V))
        V = V_new
        if delta < tol:
            return V, iteration + 1
    return V, max_iter
\end{lstlisting}

\subsection{策略迭代算法}

策略迭代交替进行策略评估和策略改进，直到策略收敛。

\textbf{算法流程}：
\begin{enumerate}[itemsep=0pt]
    \item 初始化策略$\pi_0$
    \item \textbf{策略评估}：计算$V^{\pi_k}$
    \item \textbf{策略改进}：$\pi_{k+1}(s) = \arg\max_a Q^{\pi_k}(s,a)$
    \item 如果$\pi_{k+1} = \pi_k$，停止；否则返回步骤2
\end{enumerate}

\begin{lstlisting}[caption={策略迭代算法实现}]
def run(self, max_iterations=100, eval_method='analytical'):
    policy = np.ones((n_states, n_actions)) / n_actions
    
    for iteration in range(max_iterations):
        # 策略评估
        if eval_method == 'analytical':
            V = self.evaluator.evaluate_analytical(policy)
        else:
            V, _ = self.evaluator.evaluate_iterative(policy)
        
        # 策略改进
        Q = self.mdp.compute_q_values(V)
        best_actions = np.argmax(Q, axis=1)
        new_policy = np.zeros((n_states, n_actions))
        new_policy[np.arange(n_states), best_actions] = 1.0
        
        # 检查收敛
        changed_states = np.sum(np.argmax(policy, axis=1) != best_actions)
        if changed_states == 0:
            return policy, V, iteration + 1
        policy = new_policy
\end{lstlisting}

\subsection{值迭代算法}

值迭代直接迭代值函数，不需要显式的策略评估步骤。

\textbf{贝尔曼最优方程}：
\begin{equation}
V^*(s) = \max_a \sum_{s'} P(s'|s,a)[r + \gamma V^*(s') \cdot (1-\text{done})]
\end{equation}

\begin{lstlisting}[caption={值迭代算法实现}]
def run(self, tol=1e-6, max_iterations=1000):
    V = np.zeros(n_states)
    
    for iteration in range(max_iterations):
        # 计算Q值
        Q = self.mdp.compute_q_values(V)
        # 更新值函数
        V_new = np.max(Q, axis=1)
        
        delta = np.max(np.abs(V_new - V))
        V = V_new
        
        if delta < tol:
            break
    
    # 提取最优策略
    Q = self.mdp.compute_q_values(V)
    best_actions = np.argmax(Q, axis=1)
    policy = np.zeros((n_states, n_actions))
    policy[np.arange(n_states), best_actions] = 1.0
    
    return policy, V, iteration + 1
\end{lstlisting}

\subsection{Q值计算}

Q值的正确计算是算法实现的关键。需要注意终止状态的处理：

\begin{equation}
Q(s,a) = \sum_{s'} P(s'|s,a)[r + \gamma V(s') \cdot (1-\text{done})]
\end{equation}

\begin{lstlisting}[caption={Q值计算}]
def compute_q_values(self, V):
    next_vals = V[self.T_next]          # (n, na, max_t)
    not_done  = 1.0 - self.T_done       # (n, na, max_t)
    future = self.gamma * np.sum(
        self.T_prob * next_vals * not_done, axis=2)
    return self.R_mat + future           # (n, na)
\end{lstlisting}

\section{实验结果与分析}

\subsection{算法性能比较}

\begin{table}[H]
\centering
\caption{算法性能比较}
\label{tab:performance}
\begin{tabular}{ccc}
\toprule
\textbf{指标} & \textbf{策略迭代} & \textbf{值迭代} \\
\midrule
迭代次数 & 4 & 约80 \\
运行时间(秒) & 2.35 & 1.82 \\
平均状态值 & 45.23 & 47.18 \\
最大状态值 & 89.56 & 91.32 \\
\bottomrule
\end{tabular}
\end{table}

\subsection{真实环境评估结果}

\begin{table}[H]
\centering
\caption{真实环境评估结果}
\label{tab:env_results}
\begin{tabular}{ccc}
\toprule
\textbf{策略} & \textbf{平均奖励} & \textbf{标准差} \\
\midrule
策略迭代 & 165.44 & 68.32 \\
值迭代 & 179.71 & 88.54 \\
随机策略 & 23.89 & 15.14 \\
\bottomrule
\end{tabular}
\end{table}

相比随机策略的提升：
\begin{itemize}[itemsep=0pt]
    \item 策略迭代：+592.5\%
    \item 值迭代：+652.2\%
\end{itemize}

\subsection{收敛过程可视化}

图\ref{fig:convergence}展示了两种算法的收敛过程。

\begin{figure}[H]
\centering
\includegraphics[width=0.95\textwidth]{policy_value_iteration.png}
\caption{策略迭代与值迭代收敛过程对比}
\label{fig:convergence}
\end{figure}

从图中可以观察到：
\begin{enumerate}[itemsep=0pt]
    \item 策略迭代收敛更快，仅需4轮即可收敛
    \item 值迭代需要约80次迭代才能收敛
    \item 两种算法最终收敛到相近的值函数
    \item 策略可视化显示学习到的策略符合物理直觉：杆子向右倾斜时向右施力
\end{enumerate}

\subsection{算法异同分析}

\textbf{相同点}：
\begin{enumerate}[itemsep=0pt]
    \item 都基于贝尔曼方程进行迭代计算
    \item 都能找到最优策略和最优值函数
    \item 最终收敛到相同的最优解
    \item 都需要知道环境的转移概率和奖励函数
\end{enumerate}

\textbf{不同点}：

\begin{table}[H]
\centering
\caption{策略迭代与值迭代对比}
\label{tab:comparison}
\begin{tabular}{ccc}
\toprule
\textbf{方面} & \textbf{策略迭代} & \textbf{值迭代} \\
\midrule
迭代对象 & 策略和值函数交替 & 仅迭代值函数 \\
收敛条件 & 策略稳定(不变) & 值函数收敛 \\
每次迭代 & 完整策略评估+改进 & 一次值函数更新 \\
迭代次数 & 较少(通常) & 较多 \\
每次迭代开销 & 较大(需策略评估) & 较小 \\
适用场景 & 策略空间较小 & 状态空间较大 \\
\bottomrule
\end{tabular}
\end{table}

\section{关键Bug修复说明}

在实验过程中，发现并修复了以下关键Bug：

\subsection{Bellman方程未使用done标志}

\textbf{问题描述}：终止状态的未来值未清零，导致值函数计算错误。

\textbf{正确公式}：
\begin{equation}
Q(s,a) = \sum_{s'} P(s'|s,a)[r + \gamma V(s') \cdot (1-\text{done})]
\end{equation}

\subsection{terminated与truncated混淆}

\textbf{问题描述}：采样时将terminated与truncated混淆为同一个done。

\textbf{正确理解}：
\begin{itemize}[itemsep=0pt]
    \item terminated = 真正失败（杆倒/越界），未来值应为0
    \item truncated = 步数上限，不代表失败，不应截断未来值
\end{itemize}

\subsection{策略评估收敛阈值过严}

\textbf{问题描述}：迭代策略评估容忍度$10^{-10}$过严，导致收敛极慢。

\textbf{解决方案}：改为$10^{-6}$，平衡精度和效率。

\section{结论}

本实验成功实现了策略迭代和值迭代两种动态规划算法，并将其应用于CartPole控制问题。主要结论如下：

\begin{enumerate}[itemsep=0pt]
    \item 通过状态离散化，成功将连续状态空间的CartPole问题转化为有限MDP
    \item 两种算法均能有效学习控制策略，相比随机策略提升超过500\%
    \item 策略迭代收敛更快但每次迭代开销较大，值迭代每次迭代简单但需要更多迭代次数
    \item 正确处理终止状态（done标志）是算法成功的关键
\end{enumerate}

本实验加深了对强化学习中动态规划方法的理解，为后续学习更复杂的强化学习算法奠定了基础。

\newpage

\section*{附录：完整代码}
\addcontentsline{toc}{section}{附录：完整代码}

\begin{lstlisting}[caption={策略迭代与值迭代算法完整实现（3.py）}]
import numpy as np
import matplotlib.pyplot as plt
from matplotlib import rcParams
import time

rcParams['font.sans-serif'] = ['SimHei']
rcParams['axes.unicode_minus'] = False

class CartPoleDiscreteMDP:

    def __init__(self, n_bins=8, use_env_sampling=True, n_samples=50000):
        self.n_bins = n_bins
        self.n_actions = 2
        self.action_names = ['左', '右']
        self.gamma = 0.99

        self.position_bins      = np.linspace(-2.4,  2.4,  n_bins + 1)[1:-1]
        self.velocity_bins      = np.linspace(-3.0,  3.0,  n_bins + 1)[1:-1]
        self.angle_bins         = np.linspace(-0.2,  0.2,  n_bins + 1)[1:-1]
        self.angular_velocity_bins = np.linspace(-3.0, 3.0, n_bins + 1)[1:-1]

        self.n_states = n_bins ** 4

        if use_env_sampling:
            self._build_transition_matrix_from_env(n_samples)
        else:
            self._build_transition_matrix()

        # 预计算转移张量
        self._precompute_transition_arrays()

    def discretize_state(self, state):
        cart_pos, cart_vel, pole_angle, pole_ang_vel = state
        pos_idx     = int(np.clip(np.digitize(cart_pos,     self.position_bins),      0, self.n_bins - 1))
        vel_idx     = int(np.clip(np.digitize(cart_vel,     self.velocity_bins),      0, self.n_bins - 1))
        angle_idx   = int(np.clip(np.digitize(pole_angle,   self.angle_bins),         0, self.n_bins - 1))
        ang_vel_idx = int(np.clip(np.digitize(pole_ang_vel, self.angular_velocity_bins), 0, self.n_bins - 1))
        return (pos_idx * self.n_bins ** 3 +
                vel_idx * self.n_bins ** 2 +
                angle_idx * self.n_bins +
                ang_vel_idx)

    def continuous_state(self, state_idx):
        pos_idx     = state_idx // (self.n_bins ** 3)
        remaining   = state_idx  % (self.n_bins ** 3)
        vel_idx     = remaining  // (self.n_bins ** 2)
        remaining   = remaining  % (self.n_bins ** 2)
        angle_idx   = remaining  // self.n_bins
        ang_vel_idx = remaining  %  self.n_bins

        pos_edges = np.concatenate([[-2.5], self.position_bins,         [2.5]])
        vel_edges = np.concatenate([[-3.5], self.velocity_bins,         [3.5]])
        ang_edges = np.concatenate([[-0.25], self.angle_bins,           [0.25]])
        angv_edges= np.concatenate([[-3.5], self.angular_velocity_bins, [3.5]])

        return np.array([
            (pos_edges [pos_idx]     + pos_edges [pos_idx + 1])     / 2,
            (vel_edges [vel_idx]     + vel_edges [vel_idx + 1])     / 2,
            (ang_edges [angle_idx]   + ang_edges [angle_idx + 1])   / 2,
            (angv_edges[ang_vel_idx] + angv_edges[ang_vel_idx + 1]) / 2,
        ])

    # 构建转移矩阵

    def _build_transition_matrix_from_env(self, n_samples):
        try:
            import gymnasium as gym
        except ImportError:
            import gym

        print(f"正在从真实环境采样 {n_samples} 次构建转移矩阵...")
        env = gym.make('CartPole-v1')

        # trans_data[s][a][next_s] = {'count': int, 'reward_sum': float, 'done': bool}
        trans_data = [[{} for _ in range(self.n_actions)] for _ in range(self.n_states)]

        total_samples = 0
        last_print    = 0

        while total_samples < n_samples:
            state, _ = env.reset()
            terminated = truncated = False

            while not (terminated or truncated) and total_samples < n_samples:
                s      = self.discretize_state(state)
                action = env.action_space.sample()
                next_state, reward, terminated, truncated, _ = env.step(action)
                next_s = self.discretize_state(next_state)

                # 步数超限不代表失败，未来值不应截断
                entry = trans_data[s][action]
                if next_s not in entry:
                    entry[next_s] = {'count': 0, 'reward_sum': 0.0, 'done': False}
                entry[next_s]['count']      += 1
                entry[next_s]['reward_sum'] += reward
                if terminated:
                    entry[next_s]['done'] = True  # 任何一次终止均标记

                state = next_state
                total_samples += 1
                if total_samples - last_print >= 5000:
                    print(f"  采样进度: {total_samples}/{n_samples} "
                          f"({100 * total_samples / n_samples:.0f}%)")
                    last_print = total_samples

        env.close()

        print("  正在构建转移矩阵...")
        self.P = {}
        for s in range(self.n_states):
            self.P[s] = {}
            for a in range(self.n_actions):
                total = sum(d['count'] for d in trans_data[s][a].values())
                if total == 0:
                    # 未采样到的状态：视为终止（done=True），避免错误的值传播
                    self.P[s][a] = [(1.0, s, 0.0, True)]
                else:
                    self.P[s][a] = [
                        (d['count'] / total,
                         next_s,
                         d['reward_sum'] / d['count'],
                         d['done'])
                        for next_s, d in trans_data[s][a].items()
                    ]

        print(f"采样完成，共 {total_samples} 次转移")

    def _build_transition_matrix(self):
        self.P = {}
        gravity         = 9.8
        masscart        = 1.0
        masspole        = 0.1
        total_mass      = masspole + masscart
        length          = 0.5
        polemass_length = masspole * length
        force_mag       = 10.0
        tau             = 0.02

        for s in range(self.n_states):
            self.P[s] = {}
            x, x_dot, theta, theta_dot = self.continuous_state(s)

            for a in range(self.n_actions):
                force     = force_mag if a == 1 else -force_mag
                costheta  = np.cos(theta)
                sintheta  = np.sin(theta)

                temp     = (force + polemass_length * theta_dot ** 2 * sintheta) / total_mass
                thetaacc = (gravity * sintheta - costheta * temp) / (
                    length * (4.0 / 3.0 - masspole * costheta ** 2 / total_mass))
                xacc     = temp - polemass_length * thetaacc * costheta / total_mass

                x_new        = x     + tau * x_dot
                x_dot_new    = x_dot + tau * xacc
                theta_new    = theta     + tau * theta_dot
                theta_dot_new= theta_dot + tau * thetaacc

                terminated = (x_new < -2.4 or x_new > 2.4 or
                              theta_new < -0.2095 or theta_new > 0.2095)

                if terminated:
                    # 未来值由 compute_q_values 清零
                    self.P[s][a] = [(1.0, s, 0.0, True)]
                else:
                    new_s = self.discretize_state(
                        [x_new, x_dot_new, theta_new, theta_dot_new])
                    self.P[s][a] = [(1.0, new_s, 1.0, False)]

    def _precompute_transition_arrays(self):
        n  = self.n_states
        na = self.n_actions
        max_t = max(len(self.P[s][a]) for s in range(n) for a in range(na))

        self.T_prob   = np.zeros((n, na, max_t), dtype=np.float64)
        self.T_next   = np.zeros((n, na, max_t), dtype=np.int32)
        self.T_reward = np.zeros((n, na, max_t), dtype=np.float64)
        self.T_done   = np.zeros((n, na, max_t), dtype=np.float64)

        for s in range(n):
            for a in range(na):
                for t, (prob, next_s, reward, done) in enumerate(self.P[s][a]):
                    self.T_prob  [s, a, t] = prob
                    self.T_next  [s, a, t] = int(next_s)
                    self.T_reward[s, a, t] = reward
                    self.T_done  [s, a, t] = float(done)

        # 期望即时奖励 R[s, a] = Σ_t p_t · r_t
        self.R_mat = np.sum(self.T_prob * self.T_reward, axis=2)  # (n, na)

    def compute_q_values(self, V):
        next_vals = V[self.T_next]          # (n, na, max_t)
        not_done  = 1.0 - self.T_done       # (n, na, max_t)
        future    = self.gamma * np.sum(self.T_prob * next_vals * not_done, axis=2)
        return self.R_mat + future           # (n, na)

class PolicyEvaluation:

    def __init__(self, mdp):
        self.mdp = mdp

    def evaluate_analytical(self, policy):
        n     = self.mdp.n_states
        gamma = self.mdp.gamma
        na    = self.mdp.n_actions
        max_t = self.mdp.T_prob.shape[2]

        # R_pi[s] = Σ_a π(a|s) · R_mat[s,a]
        R_pi = np.einsum('sa,sa->s', policy, self.mdp.R_mat)  # (n,)

        # P_pi[s,s'] = Σ_a π(a|s) · Σ_t p_t · (1-done_t) · δ(s', next_s_t)
        P_pi = np.zeros((n, n))
        idx  = np.arange(n)
        for a in range(na):
            for t in range(max_t):
                prob    = self.mdp.T_prob [:, a, t]          # (n,)
                next_s  = self.mdp.T_next [:, a, t]          # (n,)
                not_done= 1.0 - self.mdp.T_done[:, a, t]     # (n,)
                weight  = policy[:, a] * prob * not_done      # (n,)
                np.add.at(P_pi, (idx, next_s), weight)       # scatter-add

        A = np.eye(n) - gamma * P_pi
        try:
            V = np.linalg.solve(A, R_pi)
        except np.linalg.LinAlgError:
            V = np.linalg.lstsq(A, R_pi, rcond=None)[0]
        return V

    def evaluate_iterative(self, policy, tol=1e-6, max_iter=5000):
        V = np.zeros(self.mdp.n_states)

        for iteration in range(max_iter):
            Q     = self.mdp.compute_q_values(V)         # (n, na)
            V_new = np.einsum('sa,sa->s', policy, Q)     # (n,)  向量化！

            delta = np.max(np.abs(V_new - V))
            V     = V_new
            if delta < tol:
                return V, iteration + 1

        print(f"      警告: 策略评估未收敛, delta={delta:.2e}")
        return V, max_iter

class PolicyIteration:
    def __init__(self, mdp):
        self.mdp       = mdp
        self.evaluator = PolicyEvaluation(mdp)
        self.history   = {'iterations': [], 'policies': [], 'values': []}

    def run(self, max_iterations=100, eval_method='analytical'):
        n_states  = self.mdp.n_states
        n_actions = self.mdp.n_actions

        # 均匀随机策略初始化
        policy = np.ones((n_states, n_actions)) / n_actions
        self.history = {'iterations': [], 'policies': [], 'values': []}

        print("  开始策略迭代...")
        for iteration in range(max_iterations):
            print(f"    策略迭代第 {iteration + 1} 轮...")

            if eval_method == 'analytical':
                V = self.evaluator.evaluate_analytical(policy)
            else:
                V, eval_iters = self.evaluator.evaluate_iterative(policy)
                print(f"      策略评估迭代了 {eval_iters} 次")

            self.history['iterations'].append(iteration)
            self.history['values'].append(V.copy())
            self.history['policies'].append(policy.copy())

            # 向量化
            Q            = self.mdp.compute_q_values(V)           # (n, na)
            best_actions = np.argmax(Q, axis=1)                   # (n,)
            new_policy   = np.zeros((n_states, n_actions))
            new_policy[np.arange(n_states), best_actions] = 1.0

            old_actions     = np.argmax(policy, axis=1)
            changed_states  = np.sum(old_actions != best_actions)
            print(f"      策略变化: {changed_states} 个状态")

            policy = new_policy

            if changed_states == 0:
                print(f"  策略迭代收敛，共 {iteration + 1} 轮")
                self.history['iterations'].append(iteration + 1)
                self.history['values'].append(V.copy())
                self.history['policies'].append(policy.copy())
                return policy, V, iteration + 1

        print(f"  策略迭代达到最大轮数 {max_iterations}")
        return policy, V, max_iterations

class ValueIteration:
    def __init__(self, mdp):
        self.mdp     = mdp
        self.history = {'iterations': [], 'values': [], 'deltas': []}

    def run(self, tol=1e-6, max_iterations=1000):
        n_states  = self.mdp.n_states
        n_actions = self.mdp.n_actions
        V         = np.zeros(n_states)
        self.history = {'iterations': [], 'values': [], 'deltas': []}

        print("  开始值迭代...")
        for iteration in range(max_iterations):
            # 向量化贝尔曼最优方程
            Q     = self.mdp.compute_q_values(V)   # (n, na)
            V_new = np.max(Q, axis=1)              # (n,)  向量化！

            delta = np.max(np.abs(V_new - V))
            V     = V_new

            self.history['iterations'].append(iteration)
            self.history['values'].append(V.copy())
            self.history['deltas'].append(delta)

            if iteration % 20 == 0:
                print(f"    值迭代第 {iteration + 1} 轮, delta={delta:.2e}")

            if delta < tol:
                print(f"  值迭代收敛，共 {iteration + 1} 轮")
                break

        # 提取最优策略
        Q            = self.mdp.compute_q_values(V)
        best_actions = np.argmax(Q, axis=1)
        policy       = np.zeros((n_states, n_actions))
        policy[np.arange(n_states), best_actions] = 1.0

        return policy, V, iteration + 1

def evaluate_policy_with_env(mdp, policy, n_episodes=100):
    try:
        import gymnasium as gym
    except ImportError:
        import gym

    env           = gym.make('CartPole-v1')
    total_rewards = []

    for _ in range(n_episodes):
        state, _      = env.reset()
        episode_reward = 0
        done = truncated = False

        while not (done or truncated):
            s               = mdp.discretize_state(state)
            action          = np.argmax(policy[s])
            state, reward, done, truncated, _ = env.step(action)
            episode_reward += reward
        total_rewards.append(episode_reward)

    env.close()
    return np.mean(total_rewards), np.std(total_rewards)

def print_policy_sample(policy, mdp, n_samples=10):
    print("\n策略样本（部分状态）：")
    print("-" * 60)
    for i in range(min(n_samples, mdp.n_states)):
        state  = mdp.continuous_state(i)
        action = np.argmax(policy[i])
        print(f"状态{i:4d}: 位置={state[0]:+.2f}  角度={state[2]:+.3f} "
              f"-> 动作: {mdp.action_names[action]}")

def compare_algorithms(mdp):
    print("\n" + "=" * 60)
    print("策略迭代 vs 值迭代  比较分析")
    print("=" * 60)

    pi = PolicyIteration(mdp)
    vi = ValueIteration(mdp)

    print("\n正在运行策略迭代（解析法评估）...")
    t0 = time.time()
    pi_policy, pi_value, pi_iters = pi.run(eval_method='analytical')
    pi_time = time.time() - t0

    print("\n正在运行值迭代...")
    t0 = time.time()
    vi_policy, vi_value, vi_iters = vi.run()
    vi_time = time.time() - t0

    print("\n【算法性能比较】")
    print("-" * 50)
    print(f"{'指标':<20} {'策略迭代':<15} {'值迭代':<15}")
    print("-" * 50)
    print(f"{'迭代次数':<20} {pi_iters:<15} {vi_iters:<15}")
    print(f"{'运行时间(秒)':<20} {pi_time:<15.4f} {vi_time:<15.4f}")
    print(f"{'平均状态值':<20} {np.mean(pi_value):<15.4f} {np.mean(vi_value):<15.4f}")
    print(f"{'最大状态值':<20} {np.max(pi_value):<15.4f} {np.max(vi_value):<15.4f}")

    print("\n最优策略样本（策略迭代）")
    print_policy_sample(pi_policy, mdp, n_samples=8)

    return pi, vi, pi_policy, vi_policy

def plot_convergence(pi, vi, mdp):
    fig, axes = plt.subplots(2, 2, figsize=(14, 10))
    fig.suptitle('策略迭代 vs 值迭代  收敛分析', fontsize=16, fontweight='bold')

    # 子图1：最大状态值收敛曲线
    ax1 = axes[0, 0]
    pi_max = [np.max(v) for v in pi.history['values']]
    vi_max = [np.max(v) for v in vi.history['values']]
    ax1.plot(pi_max, 'b-o', label='策略迭代', markersize=6)
    ax1.plot(vi_max, 'r-s', label='值迭代',   markersize=4)
    ax1.set_xlabel('迭代次数', fontsize=12)
    ax1.set_ylabel('最大状态值', fontsize=12)
    ax1.set_title('收敛过程 — 最大状态值', fontsize=13)
    ax1.legend(fontsize=11); ax1.grid(True, alpha=0.3)

    # 子图2：平均状态值收敛曲线
    ax2 = axes[0, 1]
    pi_mean_v = [np.mean(v) for v in pi.history['values']]
    vi_mean_v = [np.mean(v) for v in vi.history['values']]
    ax2.plot(pi_mean_v, 'b-o', label='策略迭代', markersize=6)
    ax2.plot(vi_mean_v, 'r-s', label='值迭代',   markersize=4)
    ax2.set_xlabel('迭代次数', fontsize=12)
    ax2.set_ylabel('平均状态值', fontsize=12)
    ax2.set_title('收敛过程 — 平均状态值', fontsize=13)
    ax2.legend(fontsize=11); ax2.grid(True, alpha=0.3)

    # 子图3：值迭代 delta 对数收敛
    ax3 = axes[1, 0]
    ax3.semilogy(vi.history['deltas'], 'r-', linewidth=2, label='值迭代 δ')
    ax3.axhline(y=1e-6, color='g', linestyle='--', label='收敛阈值 1e-6')
    ax3.set_xlabel('迭代次数', fontsize=12)
    ax3.set_ylabel('δ (对数尺度)', fontsize=12)
    ax3.set_title('值迭代收敛速度', fontsize=13)
    ax3.legend(fontsize=11); ax3.grid(True, alpha=0.3)

    # 子图4：策略可视化（角度-角速度平面，固定位置/速度为中心）
    ax4 = axes[1, 1]
    nb   = mdp.n_bins
    angle_action = np.full((nb, nb), np.nan)
    final_policy = pi.history['policies'][-1]
    for i in range(nb):
        for j in range(nb):
            s = (nb // 2) * nb ** 3 + (nb // 2) * nb ** 2 + i * nb + j
            if s < mdp.n_states:
                angle_action[i, j] = np.argmax(final_policy[s])

    im = ax4.imshow(angle_action, cmap='RdYlGn', aspect='equal', vmin=0, vmax=1)
    ax4.set_title('最优策略（角度-角速度，位置/速度固定居中）', fontsize=12)
    ax4.set_xlabel('角速度离散索引', fontsize=11)
    ax4.set_ylabel('角度离散索引',   fontsize=11)
    cbar = plt.colorbar(im, ax=ax4)
    cbar.set_label('动作 (绿=右, 红=左)', fontsize=10)

    plt.tight_layout()
    save_path = 'policy_value_iteration.png'
    plt.savefig(save_path, dpi=150, bbox_inches='tight')
    plt.show()
    print(f"\n收敛图已保存至: {save_path}")

def main():
    print("\n" + "=" * 60)
    print("强化学习 — 策略迭代与值迭代算法实验")
    print("基于 CartPole-v1 离散化 MDP")
    print("=" * 60)

    # 默认使用真实环境采样
    mdp = CartPoleDiscreteMDP(n_bins=8, use_env_sampling=True, n_samples=50000)

    print("\n【MDP 环境设置】")
    print(f"  离散化 bin 数: {mdp.n_bins}")
    print(f"  总状态数:       {mdp.n_states}")
    print(f"  动作数:         {mdp.n_actions}  {mdp.action_names}")
    print(f"  折扣因子 γ:     {mdp.gamma}")
    print(f"  转移张量 max_t: {mdp.T_prob.shape[2]}")

    # ── 算法比较 ──
    pi, vi, pi_policy, vi_policy = compare_algorithms(mdp)

    # ── 可视化 ──
    plot_convergence(pi, vi, mdp)

    # ── 真实环境评估 ──
    print("\n正在用真实 CartPole 环境评估策略（100 回合）...")
    try:
        pi_mean, pi_std = evaluate_policy_with_env(mdp, pi_policy)
        vi_mean, vi_std = evaluate_policy_with_env(mdp, vi_policy)

        # 随机基线
        try:
            import gymnasium as gym
        except ImportError:
            import gym
        env            = gym.make('CartPole-v1')
        random_rewards = []
        for _ in range(100):
            state, _ = env.reset()
            total    = 0
            done = truncated = False
            while not (done or truncated):
                action = env.action_space.sample()
                state, r, done, truncated, _ = env.step(action)
                total += r
            random_rewards.append(total)
        env.close()
        rnd_mean, rnd_std = np.mean(random_rewards), np.std(random_rewards)

        print("\n【真实环境评估结果】")
        print("-" * 52)
        print(f"{'策略':<20} {'平均奖励':>10}  {'标准差':>10}")
        print("-" * 52)
        print(f"{'策略迭代':<20} {pi_mean:>10.2f}  {pi_std:>10.2f}")
        print(f"{'值迭代':<20} {vi_mean:>10.2f}  {vi_std:>10.2f}")
        print(f"{'随机策略（基线）':<20} {rnd_mean:>10.2f}  {rnd_std:>10.2f}")
        print(f"\n相比随机策略的提升：")
        print(f"  策略迭代: +{(pi_mean - rnd_mean) / rnd_mean * 100:.1f}%")
        print(f"  值  迭代: +{(vi_mean - rnd_mean) / rnd_mean * 100:.1f}%")
    except Exception as e:
        print(f"  环境评估失败（可能未安装 gymnasium）: {e}")

if __name__ == "__main__":
    main()
\end{lstlisting}

\end{document}
